% ==============================================================================
% ГЛАВА 5 (DISS-005-B2064 .. DISS-005-B2081)
% Разработка и исследование адаптивных алгоритмов БАВР с применением методов
% машинного обучения (Планируемая)
% ==============================================================================

% DISS-005-B2064
\chapter{Разработка и исследование адаптивных алгоритмов БАВР с применением методов машинного обучения (Планируемая)}
\label{ch:ml_adaptive_algorithms}

% DISS-005-B2065
\section{Введение}
\label{sec:ch05_intro}

% DISS-005-B2066
\textbf{Статус:} \textcolor{red}{В планах.}

% DISS-005-B2067
\textbf{Содержание (предполагаемое):}

% DISS-005-B2068
{\emergencystretch=1.5em Применение ML для автоматизации разметки (аннотирования) осциллограмм (Предварительный набросок уже есть):\par}

\begin{itemize}
    % DISS-005-B2069
    \item Разработка и обучение моделей (например, CNN, RNN) для классификации событий и определения их временных границ на основе датасета из Главы 3.
    % DISS-005-B2070
    \item Оценка точности и эффективности автоматической разметки.
\end{itemize}

% DISS-005-B2071

% DISS-005-B2072
Разработка адаптивных органов БАВР на основе нейронных сетей, в частности, усовершенствованного РНМ:

\begin{itemize}
    % DISS-005-B2073
    \item Формирование наборов признаков из осциллограмм (временные ряды, спектральные компоненты, статистические характеристики).
    % DISS-005-B2074
    \item Проектирование архитектур нейронных сетей, способных принимать решения о направлении мощности и блокировке/разрешении БАВР.
    % DISS-005-B2075
    \item Обучение моделей на размеченном датасете из Главы 3 (и, возможно, данных из Главы 4).
    % DISS-005-B2076
    \item Тестирование и сравнение предложенных ML-алгоритмов с традиционными и существующими адаптивными алгоритмами РНМ.
    % DISS-005-B2077
    \item Оценка вычислительной сложности и возможности реализации предложенных ML-решений на современных микропроцессорных платформах РЗА.
    % DISS-005-B2078
    \item Разработка новых адаптивных алгоритмов БАВР (в частности, РНМ) на основе методов машинного обучения, обладающих повышенной надежностью и эффективностью.
    % DISS-005-B2079
    \item Анализ применимости данных подходов для других задач РЗА
\end{itemize}

% DISS-005-B2080

% DISS-005-B2081

\clearpage
